\section{Foundations}

Mathlib organises its source code as a collection of \texttt{.lean} files arranged in a hierarchical directory structure. Each file corresponds to a \emph{module}, and its position in the directory tree is encoded by a dot-separated name such as \texttt{Mathlib.Algebra.Group}. Because every module can \texttt{import} other modules, two natural structures coexist: a tree given by the hierarchical naming, and a directed acyclic graph given by the import relation. Our goal in this section is to formalise both structures from first principles. We begin at the coarsest level---the namespace hierarchy that mirrors the file system---and then introduce the import relation that cuts across it.

\subsection{Strings and Module Names}

We adopt standard notation from formal language theory~\cite{hopcroft2006}.

\begin{definition}[Alphabet and strings]
An \emph{alphabet} $\Sigma$ is a finite non-empty set of symbols. A \emph{string} over $\Sigma$ is a finite sequence of symbols from $\Sigma$. We denote by $\Sigma^*$ the set of all strings over $\Sigma$ (the \emph{Kleene star}), and by $\Sigma^+ := \Sigma^* \setminus \{\epsilon\}$ the set of non-empty strings, where $\epsilon$ denotes the empty string.
\end{definition}

In the context of Lean and Mathlib, we work with the ASCII alphabet $\Sigma = \{0, 1, \ldots, 127\}$. A \emph{word} is a non-empty string, i.e., an element of $\Sigma^+$.

\begin{definition}[Mathlib words]
The set of \emph{Mathlib words} is
\[
  \mathcal{W} := \{w \in \Sigma^+ \mid w \text{ appears as a component in some Mathlib module name}\}.
\]
This is a finite set containing elements such as $\texttt{Mathlib}$, $\texttt{Algebra}$, $\texttt{Group}$, $\texttt{Defs}$, $\texttt{Topology}$, $\texttt{Basic}$, etc.
\end{definition}

\begin{definition}[Module names]
A \emph{module name} is a non-empty finite sequence of Mathlib words. We write $\mathcal{W}^k$ for the set of $k$-tuples of words, and define
\[
  \mathcal{W}^+ := \bigcup_{k \geq 1} \mathcal{W}^k.
\]
The set of \emph{Mathlib module names} is a finite subset
\[
  \mathcal{M} \subset \mathcal{W}^+
\]
consisting of module names that actually appear in the library.

We adopt dot notation: for a tuple $(w_1, w_2, \ldots, w_k) \in \mathcal{W}^+$, we write $w_1.w_2.\cdots.w_k$. For instance, the tuple $(\texttt{Mathlib}, \texttt{Algebra}, \texttt{Group})$ is written $\texttt{Mathlib.Algebra.Group}$.
\end{definition}

\begin{example}
Mathlib contains $|\mathcal{M}| = 7{,}563$ module names. The set $\mathcal{W}$ of Mathlib words has $|\mathcal{W}| = 3{,}855$ elements. The maximum depth (i.e., the maximum $k$ such that some module name has $k$ components) is $12$.
\end{example}

\begin{definition}[Prefix tree]
Define a partial order on $\mathcal{M}$: $m \leq m'$ if and only if $m$ is a prefix of~$m'$.
Then $(\mathcal{M}, \leq)$ is a rooted tree.
\end{definition}

\begin{example}
Consider the following subset of $\mathcal{M}$:
\begin{align*}
  \{\; & \texttt{Mathlib},\;
         \texttt{Mathlib.Algebra},\;
         \texttt{Mathlib.Algebra.Group}, \\
       & \texttt{Mathlib.Algebra.Ring},\;
         \texttt{Mathlib.Data},\;
         \texttt{Mathlib.Data.Nat} \;\}.
\end{align*}
The prefix order gives the tree:
\[
  \texttt{Mathlib}
  \begin{cases}
    \texttt{.Algebra}
    \begin{cases}
      \texttt{.Group} \\
      \texttt{.Ring}
    \end{cases} \\[6pt]
    \texttt{.Data}
    \begin{cases}
      \texttt{.Nat}
    \end{cases}
  \end{cases}
\]
Here $\texttt{Mathlib} \leq \texttt{Mathlib.Algebra} \leq \texttt{Mathlib.Algebra.Group}$,
but $\texttt{Mathlib.Algebra.Group}$ and $\texttt{Mathlib.Data.Nat}$ are incomparable.
\end{example}

\begin{definition}[File contents and import extraction]
Let $\mathrm{Content} \colon \mathcal{M} \to \Sigma^*$ assign to each module its source code. Define:
\[
  \mathrm{imports}(m) := \bigl\{m' \in \mathcal{M} \mid \texttt{"import } m'\texttt{"} \text{ is a substring of } \mathrm{Content}(m)\bigr\}.
\]
\end{definition}

\begin{example}
Suppose the file \texttt{Mathlib/Algebra/Group.lean} contains:
\begin{minted}[frame=single,fontsize=\small,breaklines,bgcolor=gray!5]{lean4}
import Mathlib.Algebra.Ring
import Mathlib.Data.Nat

class Group (G : Type) extends Monoid G where
  inv : G -> G
  mul_left_inv : forall a : G, inv a * a = 1
\end{minted}
Then $\mathrm{Content}(\texttt{Mathlib.Algebra.Group})$ is the entire string above.
Since it contains the substrings \texttt{"import Mathlib.Algebra.Ring"}
and \texttt{"import Mathlib.Data.Nat"}, we obtain
\[
  \mathrm{imports}(\texttt{Mathlib.Algebra.Group})
  = \{\texttt{Mathlib.Algebra.Ring},\; \texttt{Mathlib.Data.Nat}\}.
\]
\end{example}

\begin{definition}[Dependency graph]
The \emph{module dependency graph} of Mathlib is:
\[
  G = (\mathcal{M},\, E), \qquad E = \{(m, m') : m' \in \mathrm{imports}(m)\}.
\]
\end{definition}

\begin{example}
Continuing the previous example, suppose the import relations are:
\begin{itemize}
  \item \texttt{Mathlib.Algebra.Group} imports \texttt{Mathlib.Algebra.Ring} and \texttt{Mathlib.Data.Nat};
  \item \texttt{Mathlib.Algebra.Ring} imports \texttt{Mathlib.Data.Nat};
  \item \texttt{Mathlib.Data.Nat} imports \texttt{Mathlib.Data}.
\end{itemize}
Then the corresponding subgraph of~$G$ is:
\begin{center}
\begin{tikzpicture}[
    >=Stealth,
    node distance=1.8cm and 2.5cm,
    every node/.style={font=\small\ttfamily}
  ]
  \node (AG) {Algebra.Group};
  \node (AR) [below left=of AG] {Algebra.Ring};
  \node (DN) [below right=of AG] {Data.Nat};
  \node (D)  [below=of DN] {Data};
  \draw[->] (AG) -- (AR);
  \draw[->] (AG) -- (DN);
  \draw[->] (AR) -- (DN);
  \draw[->] (DN) -- (D);
\end{tikzpicture}
\end{center}
Each arrow $m \to m'$ indicates that $m$ imports~$m'$.
Note that the import edges cut across the prefix tree: \texttt{Algebra.Group}
depends on \texttt{Data.Nat}, which lies in a different branch.
\end{example}
